\section*{Paprika-Sahne-Haehnchen}
\ihead{}\chead{Personen:4}\ohead{}
\ifoot{Vorbereitungszeit:20 min}\cfoot{}\ofoot{Kochzeit:40 min}
{\Large Zutaten}
\begin{multicols}{2}
\begin{itemize}
\item \num{4} Hähnchenbrustfilets
\item \num{2} rote Paprikaschoten
\item \num{1} grüne Paprikaschote
\item \num{1} Zwiebel
\item \num{2} Knoblauchzehen
\item \num{1} getrocknete Chilischote
\item \num{1} Becher Sahne
\item \num{1} Becher Schmand
\item \num{1}{TL} Tomatenmark
\item \SI{125}{ml} Gemüsebrühe
\item \SI{1}{EL} Paprikapulver, edelsüß
\item \num{1}{TL} Paprikapulver, rosenscharf
\item \SI{100}{g} geriebener Käse
\item Öl
\item Salz nach Belieben
\item Paprikapulver nach Belieben
\end{itemize}
% \columnbreak
\end{multicols}
\noindent {\Large Zubereitung}\\
\\
Die Hähnchenfilets waschen und mit Küchenkrepp trocken tupfen.
Mit Salz und Paprikapulver würzen und in einer Auflaufform dicht aneinanderlegen.
Die Paprikaschoten waschen, entkernen, in schmale Streifen schneiden und auf den Filets verteilen.
Die Zwiebel in halbe Ringe schneiden und in einer Pfanne in etwas Öl andünsten.
Die Chilischote hineinzupfen und den Knoblauch pressen und hinzugeben.
Paprikapulver und Tomatenmark hinzufügen, mit der Brühe ablöschen und kurz aufkochen lassen.
Anschließend Sahne und Schmand unter die Soße rühren und mit Salz abschmecken.
Die Soße in die Auflaufform gießen, sodass Fleisch und Paprikastreifen vollständig bedeckt sind.
Den geriebenen Käse gleichmäßig darauf verteilen.
Im vorgeheizten Backofen bei \SI{180}{\celsius} Ober-/Unterhitze ca. \SI{30}{min} garen.
Dazu passen Bandnudeln oder Reis und Eisbergsalat mit Mandarinen und süß-saurer Vinaigrette.
